\section{Uniform bounds for arbitrarily many slab copies}
\label{sec:actual-slab-copy-bounds}

The centered estimates give more than bounds for a fixed number of fields.
They also control arbitrarily long strings of reflected copies with one
constant per insertion. This distinction is needed to construct bounded
operators on the reflection-positive Hilbert space. We prove that implication
on the same selected family and the same order of limits.

Write $a_k$ for its decreasing mesh scales, with $a_0=1$ and $a_k\to0$.
We use the scale inequality
\begin{equation}\label{eq:slab-mesh-gentle}
 a_{k+1}\geq c_*a_k^2\qquad(c_*>0),
\end{equation}
supplied by the selected mesh. A word is a finite ordered list of distinct
points with raw field labels; its length is denoted by $n$. Reflection
reverses the word, reflects the points in the time-zero plane, and applies
the corresponding reflected, conjugated labels. Fix a grade bound $D$ and
a bound $R$ for the norms of the raw labels. All constants below may depend
on $D,R$ and the selected family. They will not depend on $n$, the slab
width, its separation buffer, or the number of copies.

The two inputs from the selected construction are the centered service of
Section~\ref{sec:selected-centered-service}, with its uniformity over
bounded label sets, and the smooth point limits of that very family. We
retain its normalization, permutation and translation symmetries, and
reflection positivity, together with the positive-word Hilbert space and
contractive heat semigroup constructed earlier. We use the expanding order: first the fine mesh limit at a suitable fixed
period, then the cofinal period limit. The centered estimates and point
limits must use this same order. No bound at every fixed period is needed
for the argument below.

\subsection{One price for the entire reflected family}

For every $D,R$ there are $F\in\mathbb N$, $J>0$, and $P=D+1$ such that,
for $k\geq F$, every label of grade at most $D$ and norm at most $R$ has
centered root price at most
\begin{equation}\label{eq:slab-one-price}
                         J a_k^{-P}.
\end{equation}
Here the price includes every signed-coordinate transform and adjoint of
the label, as well as the chamber interiority and stencil fit needed to
apply the planted inequality.

To see the uniformity, first bound all those transformed labels by a
single $Q>0$. There are only finitely many grades and signed-coordinate
transformations. On each finite-dimensional grade space each transformation
has finite operator norm, and coefficientwise conjugation preserves the
raw norm. The adjoint of a label is the reflected conjugate
\(R_\rho\overline v\), and \(R_\rho\overline{R_sv}=R_{\rho s}\overline v\)
again belongs to this finite family, so it is covered by the same bound. We do
not assume that the signed-coordinate action is isometric. Apply the
centered service in each grade $d\leq D$, with norm bound $Q$ and exponent
slack $1$. In the expanding formulation an arbitrarily small additive
error is allowed; taking that error to be $1$ gives
$A_d a_k^{-(d+1)}+1$. Since $a_k\leq1$, this is at most
$(A_d+1)a_k^{-(d+1)}$. Take the maximum of the finitely many grade
thresholds and constants and use $d+1\leq P$. Fixed stencils fit eventually
at every such reference scale. For a fixed word only finitely many labels
and orientations occur, so their eventual mesh and period requirements
have a common tail. The price $J$ itself was selected before that word.

\subsection{Two scales for the main word and the exterior}

We first record how to combine different centered prices. At finite cutoff
let $T$ be a link-cell tiling, and let $Y$ be an aligned subdivision of it.
Suppose main insertions lie in distinct $T$-cells, and no exterior insertion
lies in those cells. The exterior insertions may share a $T$-cell, but lie
in distinct $Y$-cells. Suppose all required margins and stencil fits hold.
If the centered price of each main insertion is $A$, and that of each
exterior insertion is $X\geq0$, then
\begin{equation}\label{eq:slab-mixed-cells}
 \left|\mathbb E\left[\prod_{i\in I}f_i\prod_{j\in E}g_j\right]\right|
                 \leq A^{|I|}X^{|E|}.
\end{equation}
These are moments of the original observables; no absolute values have
been inserted inside the expectation.

Here is the extra point when several exterior fields share a coarse cell.
Move that cell to the home cell by its signed reflection, transforming all
its labels accordingly, and let $G$ be the product of its $h$ exterior
factors. Its coarse root is obtained by copying the entire product to all
$M_T$ coarse cells. The resulting $M_Th$ factors lie in distinct fine
cells: equality of their fine cells first forces equality of their coarse
cells, by alignment, and then equality of the original factors, by their
fine-scale separation. Reflections preserve the fine margins and
separation. The planted inequality in $Y$, together with the one-field
root bounds, therefore gives
\[
       |\mathbb E[T\text{-full-copy}(G)]|\leq X^{M_Th}.
\]
Taking the $M_T$th root bounds the coarse root of $G$ by $X^h$.
The empty product is covered by the same argument. The chamber-interior
property is preserved under finite products. For the main insertions the
one-field coarse root bound is the bracketing and copied-orbit argument
of Section~\ref{ex-centered-estimates}. Apply the coarse planted inequality
to these main insertions and the grouped exterior products. Their coarse
cells are distinct. Multiplying their prices proves
\eqref{eq:slab-mixed-cells}.

We also need to place all those cells without paying for the number of
main replicas. Suppose the main points have only $N_0$ residue patterns
modulo $a_m\mathbb Z^4$, and choose levels
$m\leq k\leq l\leq b$ satisfying
\begin{equation}\label{eq:slab-two-room}
 N_0(10a_k)<a_m,\qquad |E|(6a_b)<a_l.
\end{equation}
Assume main--main and main--exterior sup distances exceed $a_m$, and
exterior--exterior distances exceed $a_l$. In each coordinate, the
$N_0$ patterns forbid at most $N_0(10a_k)$ of the possible coarse offsets
when a margin of $5a_k$ is required. Thus one can first place all main
patterns in such margins. At the fine lattice level the identical
argument is a count of forbidden residues; scale divisibility gives the
same strict inequalities.

Choose the exterior offset within one aligned cell of side $a_l$.
The second inequality in \eqref{eq:slab-two-room} lets us give every
exterior point a margin of $3a_b$. This changes the main offset by at
most $a_l\leq a_k$. Allow also independent perturbations of the main
points by at most $a_k$ in every coordinate. Their original margin
$5a_k$ then leaves a margin of $3a_k$. This is why the first placement
uses $5$, rather than $3$, fine spans. The same estimate remains valid
after rounding to the fine mesh: if two coordinates differ by at most
$K a_N$ with $K$ a nonnegative integer, their ceilings in units $a_N$ differ by at most $K$ lattice
steps. Congruent unperturbed main patterns have identical coarse
residues. Thus the number of forbidden patterns stays $N_0$, even
though every occurrence may be perturbed independently.

Apply \eqref{eq:slab-mixed-cells} at that cutoff. To pass to continuum
points, first take a nearby generic configuration where the prescribed
point-limit theorem applies and where all the strict separation
inequalities remain valid. At sufficiently large period there is no
wrapping collision; at sufficiently fine mesh the rounded configuration
has the required separations and margins. The centered bounds hold on
the same mesh and period tails. Pass first to the mesh limit and then
to the period limit. Finally approach the original configuration through
generic configurations. Their density and smoothness of the limiting
correlations give \eqref{eq:slab-mixed-cells} there too. This argument
uses independent perturbations of all slots, and so does not assume
point convergence on a projection tie or on the lower-dimensional set
of exactly congruent configurations.

\subsection{Paying for the main word before choosing its exterior}

Let $w\ne\varnothing$ lie in a time slab $(0,d)$ and have buffer $r$,
where $0<r\leq1$: distinct points have sup distance at least $r$, and
every time coordinate lies in $[r,d-r]$. Choose $m\geq F$ minimal
with $a_m<r/32$, and put
\[
                   s=a_m\lfloor d/a_m\rfloor.
\]
Then $s\leq d<s+a_m$. The word remains in the slab $(0,s)$ with
buffer $r/2$, and $s>0$ because the nonempty word has both boundary
margins. Crucially, $s$ has been chosen before any exterior word.

For later use, if $j\geq F$ is the first level with $a_j<\varepsilon$,
where $0<\varepsilon\leq1$, then
\begin{equation}\label{eq:slab-first-level}
             t\varepsilon^2\leq a_j<\varepsilon,
             \qquad t=\min(a_F,c_*)>0.
\end{equation}
For $j=F$ the lower bound follows from $\varepsilon\leq1$; otherwise
$a_{j-1}\geq\varepsilon$ and \eqref{eq:slab-mesh-gentle} proves it.
Set $b_0=t/32^2$ and $u=\min(b_0/21,1)$. For $n=|w|\geq1$, choose
$k\geq F$ minimal with $a_k<\delta=u(r/n)^2$. Then
\[
 20n a_k<20n\delta=20u r^2/n\leq20u r^2<b_0r^2\leq a_m.
\]
In particular $k\geq m$. Applying \eqref{eq:slab-first-level} again gives
\begin{equation}\label{eq:slab-fine-lower}
                 a_k\geq c(r/n)^4,\qquad c=tu^2>0.
\end{equation}
All main fields and their reflected versions can therefore be paid at
this one level with price $A=Ja_k^{-P}$, and
\begin{equation}\label{eq:slab-radius-produced}
 A^n\leq C_0^n(n/r)^{4Pn},\qquad C_0=Jc^{-P}>0.
\end{equation}

To apply this to repeated insertion, write $w^\sharp=\tau_s\theta w$
and $I_wv=w\,\tau_sv$. The word
$(I_{w^\sharp}I_w)^qv$ consists of $2q$ alternating copies of the
main word followed by one translated exterior $v$. Its reflection
pairing contains $4qn$ main fields and $2|v|$ exterior fields. Every
main field is congruent modulo $a_m\mathbb Z^4$ to one of the $2n$
points of $w$ and $\theta w$, since $s$ is a multiple of $a_m$.
Translation and reflection preserve separation. Consecutive slabs
retain their buffer $r/2$, so the main points are mutually separated
and separated from the exterior by at least $r/2>a_m$.
These statements also follow directly by induction on $q$: attaching
one reflected slab and one original slab preserves the boundary
margins, adds $2n$ points, and translates the remaining word by $2s$.
The final reflected pairing doubles the main count and creates exactly
two exterior copies, independently of $q$.

Fix now an arbitrary positive exterior word $v$. Its finitely many
points have some positive mutual and boundary buffer $t_v$ and some
finite grade and norm bounds. Choose $l\geq k$ with $a_l<t_v$,
then $b\geq l$ beyond its one-field threshold and so fine that
$2|v|(6a_b)<a_l$. Let $X$ be its one-field price at level $b$.
Neither $l,b$ nor $X$ depends on $q$. The mixed estimate, with
$N_0=2n$, yields
\begin{equation}\label{eq:slab-replica-produced}
 \big\|[(I_{w^\sharp}I_w)^qv]\big\|^2
 \leq X^{2|v|} A^{4qn}
 \leq X^{2|v|}\left(C_0^n(n/r)^{4Pn}\right)^{4q}
 \qquad(q\geq0).
\end{equation}
For the empty exterior its cost is $1$ and the same proof applies.
The exterior constant may depend on $v,w,r$; its independence of $q$
is the essential conclusion.

\subsection{The actual bounded slab operators}

The repeated-insertion argument of Lemma~\ref{ex-train-bound} applies
to \eqref{eq:slab-replica-produced}, with $B_v=X^{2|v|}$ and with the
norm constant of that lemma equal to $C_0^n(n/r)^{4Pn}$. (That constant is
called $D$ there; it is unrelated to the grade bound $D$ of this section.)
It gives the operator at width $s$ with norm at most $C_0^n(n/r)^{4Pn}$. To recover the original width $d$,
first apply the heat contraction $E_{d-s}$ and then the width-$s$
insertion. On a positive word this produces
$w\,\tau_s\tau_{d-s}v=w\,\tau_dv$, exactly the required insertion.
The norm cannot increase. Thus there is a common family of slab
operators satisfying
\begin{equation}\label{eq:slab-allcopy-produced}
       \|A_d(w)\|\leq C_0^{|w|}(|w|/r)^{K|w|},\qquad K=4(D+1),
\end{equation}
for every nonempty word with the stated grade, norm and buffer bounds.
For the empty word take $A_d(\varnothing)=E_d$ with norm at most $1$;
the right side is then interpreted as $1$.

The choices used to find a bounded operator may depend on an affordable
buffer or label bound for a given word. They nevertheless produce a
single operator: two such choices agree on all planted positive words,
whose linear span is dense, and are bounded. Hence they agree everywhere.
The common operator inherits every bound obtained with any admissible
$D,R,r$. This proves the uniform all-copy estimate required by the
phase construction, on the same hierarchy as the centered service.
